\documentclass{situationist-preprint}

\title{Hypothetical Consent}
\msid{POL-02-b}
\msfamily{POL-02}
\msdiscipline{Political Philosophy}
\mstype{Companion manuscript}
\mskeywords{hypothetical consent; legitimacy; public justification; Dworkin; idealization}

\begin{document}
\maketitle

\begin{abstract}
Since almost no state has the actual consent of those it rules, many theorists appeal instead to the consent they would give under suitable conditions. Dworkin objected that a hypothetical contract is no contract at all. We argue that the objection is sound against hypothetical consent as a source of legitimacy, and that what survives it is the criterion of POL-02: a state is legitimate if it can be justified to all reasonable persons subject to it. Hypothetical consent under modest idealization is an index of justifiability; under radical idealization it is the consent of someone else. Either way, the consent does no work that the reasons do not already do.
\end{abstract}

\section{The appeal to hypothetical consent}

Consent theories of legitimacy find too little actual consent to legitimate any existing state. A natural retreat is to the consent persons \emph{would} give: Kant's original contract, which the legislator should treat as a test of whether a law could have come from the united will of the people \citep{Kant1793}, and the contractualist traditions that descend from it. The appeal promises to keep what was attractive in consent, the idea that legitimate authority rests on the will of those subject to it, while dropping what made it unrealistic.

\citet{Dworkin1973} replied that a hypothetical contract is not a pale form of an actual contract but no contract at all. That I would have agreed to sell you my painting for a hundred dollars yesterday, had you asked, does not oblige me to sell it to you today for that price. If the hypothetical agreement has any force, it is the force of the reasons that would have led me to make it.

\section{The dilemma}

Let a person's hypothetical consent to the state's authority be the consent she would give under some idealization $I$ of her beliefs, values and circumstances.

\begin{proposition}\label{p:dilemma}
Either $I$ is modest, and then hypothetical consent coincides with justifiability in the sense of Definition~3.3 of POL-02-a; or $I$ is radical, and then hypothetical consent is not consent of the person subject to the state.
\end{proposition}

\begin{proof}
Suppose $I$ corrects only factual errors and inferential mistakes and leaves the person's values and commitments as they are. Then she would consent under $I$ just in case there is a justification of the state's authority that proceeds from considerations she has sufficient reason to accept from her own standpoint, so corrected; that is Definition~3.3 of POL-02-a. Suppose instead that $I$ changes her values or commitments. Then the person who would consent is one whose standpoint differs from hers, and the consent is that person's, not hers.
\end{proof}

On the first horn, the counterfactual act adds nothing: what makes the state legitimate with respect to her is the availability of a justification she has sufficient reason to accept, and the claim that she would consent is a way of saying so. On the second horn, the appeal violates the requirement that motivates consent theory in the first place, that the authority be grounded in the standpoint of the person it binds. \citet{Stark2000} reaches a similar conclusion: hypothetical consent is better understood as a device of justification than as a source of obligation.

This vindicates the structure of POL-02-a. Its criterion does not claim that the state binds through consent, actual or hypothetical. It holds that coercion of a reasonable person is permissible only if justifiable to her (Principle~P1 of POL-02-a) and is permissible if it is (Principle~P2), and it explains why actual consent, where present, also suffices (Proposition~5.3 of POL-02-a): actual consent is one way, not the only way, in which a justification can be available to a person.

\section{Objections}

\begin{objection}
Hypothetical consent matters because it respects autonomy: to act on someone as she would have consented is to treat her as a self-governing agent, even when no actual consent was given, as in emergency medicine.
\end{objection}

\begin{reply}
The medical cases support our view. The surgeon who operates on an unconscious patient is guided by what this patient, with her own values, would accept; that is the modest horn, and what licenses the operation is the patient's reasons. Where autonomy-based accounts give hypothetical consent independent weight, as \citet{Enoch2017} considers, the weight is greatest where actual consent is impossible, and political subjection is not such a case. The state that could ask, and does not, cannot invoke what the subject would have said.
\end{reply}

\begin{objection}
If hypothetical consent reduces to justifiability, the criterion of POL-02 inherits the indeterminacy of idealization: different idealizations yield different verdicts.
\end{objection}

\begin{reply}
The criterion fixes the idealization as modest, and so does not range over idealizations. The residual vagueness of \enquote{modest} is a vagueness in the concept of a person's own reasons, which any account of respect for persons must accept.
\end{reply}

\section{Conclusion}

Hypothetical consent cannot legitimate a state by itself. Under modest idealization it is a restatement of justifiability to the person; under radical idealization it is not her consent. The criterion of POL-02, justifiability to all reasonable persons subject to the state, keeps what was right in the appeal to consent and drops what was wrong.

\begin{thebibliography}{9}
\bibitem[Dworkin(1973)]{Dworkin1973} Dworkin, R. (1973). The original position. \emph{University of Chicago Law Review}, 40(3), 500--533.
\bibitem[Enoch(2017)]{Enoch2017} Enoch, D. (2017). Hypothetical consent and the value(s) of autonomy. \emph{Ethics}, 128(1), 6--36.
\bibitem[Kant(1793)]{Kant1793} Kant, I. (1793). On the common saying: \enquote{This may be true in theory, but it does not apply in practice}. In \emph{Political Writings}, ed. H.~S. Reiss, trans. H.~B. Nisbet, 2nd edn. Cambridge University Press, 1991.
\bibitem[Stark(2000)]{Stark2000} Stark, C.~A. (2000). Hypothetical consent and justification. \emph{The Journal of Philosophy}, 97(6), 313--334.
\end{thebibliography}

\end{document}
